\section{与设计的差异、未覆盖项与反思}

\subsection{与实验二设计的五处差异}

任务书允许"实现与设计不一致"，但要求说明原因。本次一共有五处差异（\tabref{tab:diff}），其中第 3 处需要回改设计文档，其余是实现层面的调整。

\begin{table}[htbp]
\centering
\caption{实验二设计到实验三实现的差异}\label{tab:diff}
\zihao{-5}
\renewcommand{\arraystretch}{1.16}
\begin{tabularx}{\textwidth}{@{}L{0.8cm}L{4.2cm}L{3.6cm}Y@{}}
\toprule
序号 & 设计 & 实现 & 原因 \\
\midrule
1 & \code{Item}、\code{Reservation} 自带业务方法 & 方法上移到 \code{ItemService}、\code{TradeService} & 方法需要跨聚合读写，实体会被迫持有仓储，事务边界不清（详见 2.3） \\
2 & 查询为独立组件 M4 & 并入商品服务 & 查询条件就是商品属性，拆开会造成两个服务同时读 \code{Item} \\
3 & \code{ItemRevision} 只暂存标题、价格、图片 & 增加 \code{categoryId} & 改分类也属于"修改商品信息"，不带分类会导致审核后内容与分类不一致 \\
4 & 新增 \code{ServiceBase}、\code{Store}、\code{Table} & 三个实现类 & 收敛"令牌换用户""落盘""泛型查询"三处重复代码 \\
5 & 前后端分离的 B/S 结构，用户端 H5 & 单机命令行界面 & 本次重点是规则正确性；服务层与界面层已隔离，换界面不动业务 \\
\bottomrule
\end{tabularx}
\end{table}

第 3 处建议回改实验二类图，已在仓库的类图（\code{UML/UML类图.puml}）中体现；第 5 处属于形态调整，实验二报告第 1 节写的"形态为前后端分离的 B/S 结构"在本次并未实现，这一点在下一小节作为未覆盖项明确列出。

\subsection{未覆盖或简化的部分}

诚实记录四条，避免把"没做"说成"做了"：

\begin{enumerate}[leftmargin=1.6em,itemsep=2pt,topsep=3pt]
\item \textbf{REQ-011（查询响应时间）未做实测。}需求约定"1 万条商品、20 并发、100 次查询中至少 95 次在 1 秒内"。当前实现是单机内存表 + 整库快照写盘，20 并发的语义与 Web 服务不同，测出来的数字不能直接对照这条非功能需求。合理的做法是在实验四的性能测试里针对 \code{ItemService.search()} 做基准，用生成的 1 万条数据测单次查询耗时，并说明并发模型差异。
\item \textbf{REQ-012 的 HTTPS 未实现。}桌面程序没有 HTTP 通信，这条需求在本形态下不适用；需求中"密码加密存储、校验不还原明文"的部分已实现（PBKDF2-HMAC-SHA256，12 万轮，随机盐，密文比对）。
\item \textbf{REQ-009 仍为待确认。}实现上按"管理员可维护分类"做了增删改查，并加了"被商品引用的分类不能删除""分类名不能重复"两条约束。这条需求在实验一就是待确认状态，实现只是给了一个可用的默认方案，等确认后再调整。
\item \textbf{没有图形界面。}实验二的 10 张界面稿没有落地。控制台界面能覆盖全部用例，但演示效果差，这是本次最明显的短板。
\end{enumerate}

\subsection{反思}

\textbf{第一，"能跑"离"对"还有很远。}四个真实缺陷里有三个（审核期间被预约、完成只通知卖家、必填校验过早）在功能演示时都不报错，是演示脚本和测试里的"反例"把它们撞出来的。如果只按正常流程点一遍，这些缺陷会一直留在代码里。以后再写这类有规则的模块，我会先列"应该被拒绝的场景清单"，再写实现。

\textbf{第二，设计的价值在写代码时才显现。}实验二的状态图让我在写 \code{state.py} 时只需要翻译 9 条箭头，不用边写边想"已预约能不能下架"；而时序图里"幂等键查重放在状态校验之前"这个顺序，直接决定了重发请求不会被误判为冲突。反过来，设计里没画清楚的地方（改分类要不要一起待审）就变成了实现期的返工。

\textbf{第三，静态检查应该早跑、常跑。}第一次跑 pylint 是写完 3000 行之后，27 条告警里有 22 条是行宽，属于"写的时候就该注意"的问题。按模块跑、每轮清干净，改起来是几秒钟的事；攒到最后再改，就是一次几百行的机械劳动，还容易改错。

\textbf{下一步。}如果继续做，优先补三件事：给 \code{ItemService.search()} 做性能基准以满足 REQ-011；把命令行换成 Tkinter 界面复用实验二的界面稿；把 \code{Store} 换成 SQLite，用唯一索引在执行层再保证一次"同一商品至多一条有效预约"，而不是只依赖服务层的锁。
